% from ega1-1-fr.tex
\oldpage[I]{97}
Posons $A=A(X)$, $B=A(Y)$. La condition est nécessaire, car on a vu
(\hyperref[I.1.6.1-fr]{1.6.1}) que $\widetilde\varphi_x^{\#}$ est
l'homomorphisme de $B_{{}^a\varphi(x)}$ dans $A_x$ déduit canoniquement de
$\varphi$, et par définition de
${}^a\varphi(x)=\varphi^{-1}(\mathfrak{j}_x)$, cet homomorphisme est local.

Prouvons que la condition est suffisante. Par définition, $\theta$ est un
homomorphisme $\mathscr{O}_Y\to\psi_*(\mathscr{O}_X)$, et on en déduit
canoniquement un homomorphisme d'anneaux
\[
  \varphi=\Gamma(\theta):
  B=\Gamma(Y,\mathscr{O}_Y)
  \longrightarrow
  \Gamma(Y,\psi_*(\mathscr{O}_X))=\Gamma(X,\mathscr{O}_X)=A.
\]

L'hypothèse sur $\theta_x^{\#}$ permet de déduire de cet homomorphisme, par
passage aux quotients, un monomorphisme $\theta^x$ du corps des restes
$k(\psi(x))$ dans le corps des restes $k(x)$, tel que, pour toute section
$f\in\Gamma(Y,\mathscr{O}_Y)=B$, on ait
$\theta^x(f(\psi(x)))=\varphi(f)(x)$. La relation $f(\psi(x))=0$ est donc
équivalente à $\varphi(f)(x)=0$, ce qui signifie que
$\mathfrak{j}_{\psi(x)}=\mathfrak{j}_{{}^a\varphi(x)}$, et s'écrit encore
$\psi(x)={}^a\varphi(x)$ pour tout $x\in X$, ou $\psi={}^a\varphi$. On sait
aussi que le diagramme
\phantomsection\label{I.1.7.3.diagram-fr}
\[
  \xymatrix{
    B=\Gamma(Y,\mathscr{O}_Y)\ar[r]^\varphi\ar[d] &
    \Gamma(X,\mathscr{O}_X)=A\ar[d]\\
    B_{\psi(x)}\ar[r]_{\theta_x^{\#}} & A_x
  }
\]
est commutatif (\hyperref[0.3.7.2-fr]{0, 3.7.2}), ce qui signifie que
$\theta_x^{\#}$ est égal à l'homomorphisme
$\varphi_x:B_{\psi(x)}\to A_x$ déduit canoniquement de $\varphi$
(\hyperref[0.1.5.1-fr]{0, 1.5.1}). Comme la donnée des $\theta_x^{\#}$
caractérise complètement $\theta^{\#}$, et par suite aussi $\theta$
(\hyperref[0.3.7.1-fr]{0, 3.7.1}), on en conclut que l'on a
$\theta=\widetilde\varphi$, par définition de $\widetilde\varphi$
(\hyperref[I.1.6.1-fr]{1.6.1}).

Nous dirons qu'un morphisme $(\psi,\theta)$ d'espaces annelés satisfaisant à
la condition de (\hyperref[I.1.7.3-fr]{1.7.3}) est un \emph{morphisme de
schémas affines}.

\begin{corollary}[1.7.4]
\label{I.1.7.4-fr}
Si $(X,\mathscr{O}_X)$, $(Y,\mathscr{O}_Y)$ sont deux schémas affines, il
existe un isomorphisme canonique de l'ensemble de morphismes de schémas
affines
$\operatorname{Hom}((X,\mathscr{O}_X),(Y,\mathscr{O}_Y))$ sur l'ensemble
d'homomorphismes d'anneaux de $B$ dans $A$, où
$A=\Gamma(\mathscr{O}_X)$ et $B=\Gamma(\mathscr{O}_Y)$.
\end{corollary}

On peut encore dire que les foncteurs $(\operatorname{Spec}(A),\widetilde A)$
en $A$ et $\Gamma(X,\mathscr{O}_X)$ en $(X,\mathscr{O}_X)$ définissent une
\emph{équivalence} de la catégorie des anneaux commutatifs et de la catégorie
duale de la catégorie des schémas affines (T, I, 1.2).

\begin{corollary}[1.7.5]
\label{I.1.7.5-fr}
Si $\varphi:B\to A$ est surjectif, le morphisme correspondant
$({}^a\varphi,\widetilde\varphi)$ est un monomorphisme d'espaces annelés
\emph{(cf. (\hyperref[I.4.1.7-fr]{4.1.7}))}.
\end{corollary}

On sait en effet que ${}^a\varphi$ est injective
(\hyperref[I.1.2.5-fr]{1.2.5}), et comme $\varphi$ est surjectif, pour tout
$x\in X$, $\varphi_x^{\#}:B_{{}^a\varphi(x)}\to A_x$, qui se déduit de
$\varphi$ par passage aux anneaux de fractions, est aussi surjectif
(\hyperref[0.1.5.1-fr]{0, 1.5.1})~; d'où la conclusion
(\hyperref[0.4.1.1-fr]{0, 4.1.1}).
